\scanpage{43}
\diagram{\scriptsize\resizebox{\linewidth}{!}{\(\begin{array}{r@{\quad}l@{\quad}r@{}l}
4&\mathcal A_4&
\begin{array}{r}3\\1\end{array}&
\left|\begin{array}{l}
(3,3',2)\ (3',2,3)\ (2,3,3')\quad[(\deleted{\unsure{3}}2,3',3)\ (3',3,2)\ (3,2,3')]\ \deleted{\text{standart}}\\ 
(3,3,3)\qquad[(3',3',3')]
\end{array}\right.\\[.8em]
9&\unsure{\mathfrak S_4|\mathscr S_4}&
\begin{array}{r}6\\3\end{array}&
\left|\begin{array}{ll}
(4,3,2)\ (3,2,4)\ (2,4,3)\ (2,3,4)\ (3,4,2)\ (4,2,3)&\text{standart}\\ 
(4,4,3)\ (4,3,4)\ (3,4,4)&
\end{array}\right.\\[.8em]
32&\mathcal A_5&
\begin{array}{r}6\\6\end{array}&
\left\{\begin{array}{ll}
(5,3,2)\ (3,2,5)\ (2,5,3)\ (2,3,5)\ (3,5,2)\ (5,2,3)&\text{standart }\alpha\\ 
(5',3,2)\ \deleted{(3,2,5')}\ (2,5',3)\ (2,3,5')\ (3,5',2)\ (5',2,3)&\text{id}\quad\alpha'
\end{array}\right.\\[.3em]
&&\begin{array}{r}3\\3\end{array}&
\left\{\begin{array}{l}
(5,5,3)\ (5,3,5)\ (3,5,5)\\ 
(5',5',3)\ (5',3,5')\ (3,5',5')
\end{array}\right.\\[.3em]
&&6&(5,5',3)\ (5',3,5)\ (3,5,5')\ (3,5',5)\ (5',5,3)\ (5,3,5')\\[.3em]
&&6&(5,5',2)\ (5',2,5)\ (2,5,5')\ (2,5',5)\ (5',5,2)\ (5,2,5')\\[.3em]
&&\begin{array}{r}1\\1\end{array}&
\left\{\begin{array}{l}(5,5,5)\\(5',5',5')\end{array}\right.\\ 
\overline{45}&&&
\end{array}\)}}

\margin{(ou 16 si on regarde mod action de $\mathfrak S_5$)}

Question : un syst. de Schwartz d'un
groupe fini est-il connu \inserted{mod conjugaison,} quand on
connaît chacun de ses trois comp.
mod conjugaison ? OK dans $\mathcal A_4,\mathfrak S_4,\mathcal A_5$
-- regarder groupes diédraux, et \uline{$\mathfrak S_5$\ldots}

% archive sheet number: 42
